\documentclass[aspectratio=169]{beamer}

\usetheme{Madrid}
\usecolortheme{default}

\usepackage{amsmath,amssymb,amsfonts,bm}
\usepackage{physics}
\usepackage{mathtools}
\usepackage{booktabs}

\title{Stochastic Reconfiguration in Quantum Monte Carlo}
\author{Morten Hjorth-Jensen}
\date{FYS4411/9411, May 8, 2026}

\newcommand{\Eloc}{E_{\mathrm{loc}}}
\newcommand{\dkl}{D_{\mathrm{KL}}}

\begin{document}

\frame{\titlepage}

% ══════════════════════════════════════════════════════════════════════════════
%  CONTENT
% ══════════════════════════════════════════════════════════════════════════════
\begin{frame}{Content}
\begin{columns}[T]
\begin{column}{0.48\textwidth}
\textbf{Stochastic reconfiguration}
\begin{itemize}
  \item imaginary-time evolution and variational projection
  \item tangent-space expansion of the trial state
  \item derivation of \(S\,\delta\bm{\alpha}=-\delta\tau\,\bm{g}\)
  \item logarithmic derivatives, metric matrix, force vector
  \item regularization and algorithmic summary
\end{itemize}
\end{column}
\begin{column}{0.48\textwidth}
\textbf{The score function and Fisher information}
\begin{itemize}
  \item why we use \(\log p(x)\)
  \item likelihood, log-likelihood, and the score
  \item relative sensitivity and information content
  \item entropy, additivity, and information geometry
  \item Fisher information matrix; quantum FIM \(F_{ij}=4S_{ij}\)
  \item natural gradient; three-way equivalence with SR
  \item Cram\'er--Rao bound and regularization
\end{itemize}
\end{column}
\end{columns}
\end{frame}

% ══════════════════════════════════════════════════════════════════════════════
%  PART I  –  STOCHASTIC RECONFIGURATION
% ══════════════════════════════════════════════════════════════════════════════

\begin{frame}{Setup and notation}
We consider a variational wave function
\[
\ket{\Psi(\bm{\alpha})},
\qquad
\bm{\alpha}=(\alpha_1,\alpha_2,\dots,\alpha_p),
\]
depending on \(p\) real parameters.

\medskip

The normalized state is
\[
\ket{\overline{\Psi}(\bm{\alpha})}
=
\frac{\ket{\Psi(\bm{\alpha})}}
     {\sqrt{\braket{\Psi(\bm{\alpha})}{\Psi(\bm{\alpha})}}}.
\]

We wish to optimize the parameters such that the energy
\[
E(\bm{\alpha})
=
\frac{\bra{\Psi(\bm{\alpha})}H\ket{\Psi(\bm{\alpha})}}
     {\braket{\Psi(\bm{\alpha})}{\Psi(\bm{\alpha})}}
\]
is minimized.
\end{frame}

% ─────────────────────────────────────────────────────────────────────────────
\begin{frame}{Imaginary-time evolution as a route to the ground state}
The exact imaginary-time evolution is
\[
\ket{\Psi(\tau)}=e^{-\tau H}\ket{\Psi(0)}.
\]

Expanding in energy eigenstates \(H\ket{n}=E_n\ket{n}\):
\[
\ket{\Psi(\tau)}=\sum_n c_n\,e^{-\tau E_n}\ket{n}.
\]

If \(c_0\neq 0\), then for large \(\tau\),
\[
\ket{\Psi(\tau)} \propto e^{-\tau E_0}\ket{0},
\]
so imaginary-time propagation filters out the ground state.
\end{frame}

% ─────────────────────────────────────────────────────────────────────────────
\begin{frame}{Why a variational projection is needed}
In Variational Monte Carlo we restrict to the variational manifold
\[
\mathcal{M}
=\left\{\ket{\Psi(\bm{\alpha})}\;\middle|\;\bm{\alpha}\in\mathbb{R}^p\right\}.
\]

The exact state
\(
(1-\delta\tau H)\ket{\Psi(\bm{\alpha})}
\)
will in general \emph{not} belong to \(\mathcal{M}\).

\medskip

Idea of stochastic reconfiguration:
\begin{itemize}
  \item apply one infinitesimal imaginary-time step,
  \item project the resulting state back onto \(\mathcal{M}\),
  \item interpret the projection as a parameter update
        \(\bm{\alpha}\to\bm{\alpha}+\delta\bm{\alpha}\).
\end{itemize}
\end{frame}

% ─────────────────────────────────────────────────────────────────────────────
\begin{frame}{First-order imaginary-time step}
For a small time step \(\delta\tau\),
\(
e^{-\delta\tau H}=1-\delta\tau H+\mathcal{O}(\delta\tau^2).
\)

A convenient norm-preserving first-order approximation is
\[
\ket{\Psi_{\mathrm{exact}}'}
\approx
\left[1-\delta\tau\bigl(H-E(\bm{\alpha})\bigr)\right]\ket{\Psi(\bm{\alpha})}.
\]

Since \(\bra{\Psi}(H-E)\ket{\Psi}=0\), the vector \((H-E)\ket{\Psi}\)
is orthogonal to \(\ket{\Psi}\) to first order, so subtracting
\(E(\bm{\alpha})\) removes the trivial norm-rescaling component and
moves the state tangentially in projective Hilbert space.
\end{frame}

% ─────────────────────────────────────────────────────────────────────────────
\begin{frame}{Tangent-space expansion and projection}
The variational state at shifted parameters is
\[
\ket{\Psi(\bm{\alpha}+\delta\bm{\alpha})}
=\ket{\Psi}+\sum_k\delta\alpha_k\ket{\Psi_k}+\mathcal{O}(\delta\alpha^2),
\qquad
\ket{\Psi_k}\equiv\pdv{}{\alpha_k}\ket{\Psi}.
\]

The \textbf{mismatch vector}
\(\ket{\Delta}=\ket{\Psi_{\mathrm{var}}'}-\ket{\Psi_{\mathrm{exact}}'}\)
has to first order
\[
\ket{\Delta}
=\sum_k\delta\alpha_k\ket{\Psi_k}+\delta\tau(H-E)\ket{\Psi}.
\]

Minimizing \(\mathcal{D}=\braket{\Delta}{\Delta}\) and differentiating
with respect to \(\delta\alpha_i\):
\[
\sum_j\braket{\Psi_i}{\Psi_j}\,\delta\alpha_j
=-\delta\tau\,\Re\bra{\Psi_i}(H-E)\ket{\Psi}.
\]
\end{frame}

% ─────────────────────────────────────────────────────────────────────────────
\begin{frame}{Logarithmic derivatives and the SR matrix}
In configuration space define the \textbf{logarithmic derivative}:
\[
O_k(x)
=\frac{1}{\Psi(x)}\pdv{\Psi(x)}{\alpha_k}
=\pdv{\ln\Psi(x)}{\alpha_k}.
\]

Monte Carlo expectation values:
\(
\langle A\rangle=N_s^{-1}\sum_\ell A(x^{(\ell)})
\),
\;
\(
P(x)=|\Psi(x)|^2/\braket{\Psi}{\Psi}
\).

\medskip

Working with normalized states, for real \(\Psi\):
\[
S_{ij}
=\langle O_i O_j\rangle-\langle O_i\rangle\langle O_j\rangle
\quad\text{(covariance of log-derivatives)},
\]
\[
g_i
=\langle O_i\Eloc\rangle-\langle O_i\rangle\langle\Eloc\rangle
\quad\text{(covariance with local energy)},
\]
where \(\Eloc(x)=(H\Psi)(x)/\Psi(x)\).
\end{frame}

% ─────────────────────────────────────────────────────────────────────────────
\begin{frame}{Final SR linear system}
Collecting everything:
\[
\boxed{
S\,\delta\bm{\alpha}=-\delta\tau\,\bm{g},
\qquad
\delta\bm{\alpha}=-\delta\tau\,S^{-1}\bm{g},
}
\]
\[
S_{ij}=\langle O_iO_j\rangle-\langle O_i\rangle\langle O_j\rangle,
\qquad
g_i=\langle O_i\Eloc\rangle-\langle O_i\rangle\langle\Eloc\rangle.
\]

\medskip

\textbf{Comparison with gradient descent:}
ordinary GD uses \(\delta\bm{\alpha}_{\mathrm{GD}}=-\eta\,\bm{g}\);
SR uses \(\delta\bm{\alpha}_{\mathrm{SR}}=-\eta\,S^{-1}\bm{g}\).
The inverse metric rescales and mixes parameter directions according
to the geometry of the wave-function manifold, giving better conditioning
and a direct connection to imaginary-time evolution.
\end{frame}

% ══════════════════════════════════════════════════════════════════════════════
%  PART II  –  THE SCORE FUNCTION AND FISHER INFORMATION
% ══════════════════════════════════════════════════════════════════════════════

\begin{frame}{}
\vfill
\begin{center}
{\Large\bfseries Part~II}\\[10pt]
{\large The Score Function, Fisher Information,\\[4pt]
and Their Equivalence with Stochastic Reconfiguration}
\end{center}
\vfill
\end{frame}

% ── Why log p? ────────────────────────────────────────────────────────────────

\begin{frame}{Why does \(\log p(x)\) appear everywhere?}
In statistics, machine learning, information theory, and statistical
physics, the quantity \(\log p(x)\) is ubiquitous.

\medskip

\begin{columns}[T]
\begin{column}{0.48\textwidth}
\textbf{Where it appears}
\begin{itemize}
  \item Maximum likelihood estimation
  \item Fisher information
  \item Cross-entropy loss
  \item Shannon entropy
  \item Bayesian inference
  \item Statistical mechanics
  \item Variational inference
\end{itemize}
\end{column}
\begin{column}{0.48\textwidth}
\textbf{One-sentence answer}

Probabilities \emph{multiply} while optimisation prefers
\emph{additive} structures.
The logarithm converts products into sums and relative
changes into absolute ones.

\medskip

The following slides make this precise.
\end{column}
\end{columns}
\end{frame}

% ─────────────────────────────────────────────────────────────────────────────
\begin{frame}{Likelihood and why we take its logarithm}
Observe independent data \(x_1,\dots,x_N\) drawn from \(p(x|\theta)\).
The \textbf{likelihood} is
\[
P(D|\theta)=\prod_{i=1}^N p(x_i|\theta).
\]

Products are inconvenient: numerically unstable, exponentially small,
hard to differentiate.  Taking the logarithm:
\[
\log P(D|\theta)=\sum_{i=1}^N \log p(x_i|\theta).
\]

Because \(\log\) is monotonic,
\[
\arg\max_\theta P(D|\theta)
=\arg\max_\theta \log P(D|\theta),
\]
so maximising the likelihood and the log-likelihood are \emph{equivalent}.
The log-likelihood is far easier to work with analytically and numerically.
\end{frame}

% ── Score function ────────────────────────────────────────────────────────────

\begin{frame}{The score function}
The \textbf{score function} is the gradient of the log-likelihood:
\[
s(\theta)
=\partial_\theta \log p(x|\theta)
=\frac{1}{p(x|\theta)}\,\partial_\theta\, p(x|\theta).
\]

The second form shows that the score measures the
\textbf{relative sensitivity} of the distribution to the parameter:
\[
s(\theta)=\frac{\delta p\,/\,\partial\theta}{p}.
\]

\medskip

\textbf{Why relative, not absolute?}
Suppose \(p(x)\ll 1\).  A small absolute change in \(p\) may
correspond to a large relative change — and it is the relative
change that matters for distinguishability.
The logarithmic derivative \(\delta p/p\) captures the natural
geometry of probability distributions.
\end{frame}

% ─────────────────────────────────────────────────────────────────────────────
\begin{frame}{Key identities of the score function}
\textbf{Identity 1 — zero mean.}
Differentiating the normalisation \(\int p(x|\theta)\,dx=1\):
\[
\mathbb{E}[\partial_\theta \log p]
=\int p\,\frac{\partial_\theta p}{p}\,dx
=\partial_\theta\!\int p\,dx=0.
\]

\textbf{Identity 2 — Hessian form of Fisher information.}
Differentiating the zero-mean identity again:
\[
\mathbb{E}\bigl[(\partial_\theta \log p)^2\bigr]
=-\mathbb{E}\bigl[\partial_\theta^2 \log p\bigr].
\]

\medskip

These two identities hold for any smooth parametric family and rely
crucially on the logarithm.  They are the foundation of both maximum
likelihood theory and the Fisher information matrix.
\end{frame}

% ── Information theory ────────────────────────────────────────────────────────

\begin{frame}{Information content and entropy}
Define the \textbf{information content} of an event \(x\):
\[
I(x)=-\log p(x).
\]

Rare events (\(p(x)\ll 1\)) carry large information;
common events carry little.

\medskip

For independent events \(A\) and \(B\), probabilities multiply
but information \emph{adds}:
\[
p(A,B)=p(A)\,p(B)
\quad\Longrightarrow\quad
I(A,B)=I(A)+I(B).
\]

The \textbf{Shannon entropy} is the expected information content:
\[
S=-\sum_x p(x)\log p(x)=\mathbb{E}[I(x)].
\]

Entropy measures the average surprise of the distribution.
The logarithm is the \emph{unique} function (up to a constant)
that makes information additive for independent events.
\end{frame}

% ─────────────────────────────────────────────────────────────────────────────
\begin{frame}{Connection to statistical mechanics and machine learning}
\begin{columns}[T]
\begin{column}{0.48\textwidth}
\textbf{Statistical mechanics}\\[4pt]
The Boltzmann distribution is
\[
p(x)=\frac{e^{-\beta E(x)}}{Z}.
\]
Taking logarithms:
\[
\log p(x)=-\beta E(x)-\log Z,
\]
so energy is \emph{proportional} to log-probability.
This makes the connection to thermodynamics exact
and allows free-energy calculations via log-likelihoods.
\end{column}
\begin{column}{0.48\textwidth}
\textbf{Machine learning}\\[4pt]
Many generative models are built on log-probabilities:
\begin{itemize}
  \item cross-entropy loss:
        \(L=-\sum_i y_i\log p_i\)
  \item variational autoencoders
  \item diffusion models
  \item Boltzmann machines
  \item Bayesian neural networks
\end{itemize}
In each case the logarithm converts the product structure of
joint probabilities into an additive, differentiable objective.
\end{column}
\end{columns}
\end{frame}

% ── Fisher information matrix ─────────────────────────────────────────────────

\begin{frame}{Fisher information matrix — definition}
Let \(p(x;\bm{\alpha})\) be a parametric family of distributions
and \(s_k=\partial_{\alpha_k}\log p(x;\bm{\alpha})\) the score.
Because \(\mathbb{E}_p[s_k]=0\), the \textbf{Fisher information
matrix} (FIM) equals the covariance of the score:
\[
\boxed{
F_{ij}
=\mathbb{E}_p[s_i\,s_j]
=\mathrm{Cov}_p(s_i,s_j)
=-\mathbb{E}_p\!\left[
  \frac{\partial^2\log p}{\partial\alpha_i\,\partial\alpha_j}
\right].}
\]

\medskip

Interpretations:
\begin{itemize}
  \item \textbf{Curvature of the log-likelihood:} large \(F_{ii}\)
        means the likelihood surface is sharply curved around
        the parameter \(\alpha_i\).
  \item \textbf{Distinguishability:} large \(F_{ii}\) means nearby
        distributions are easy to tell apart.
  \item \textbf{Riemannian metric:} \(F\) defines distances between
        distributions via \(ds^2=d\bm{\alpha}^\top F\,d\bm{\alpha}\).
\end{itemize}
\end{frame}

% ─────────────────────────────────────────────────────────────────────────────
\begin{frame}{Fisher information as the curvature of KL divergence}
\textbf{Theorem.}  For \(\bm{\alpha}'=\bm{\alpha}+d\bm{\alpha}\),
\[
\dkl\!\bigl(p(\cdot;\bm{\alpha})\,\|\,p(\cdot;\bm{\alpha}')\bigr)
=\frac{1}{2}\,d\bm{\alpha}^\top F(\bm{\alpha})\,d\bm{\alpha}
+\mathcal{O}(d\alpha^3).
\]

\textbf{Proof sketch.}  Taylor-expand the KL divergence to
second order in \(d\bm{\alpha}\):
\begin{multline*}
\dkl(p\|p')
=-\mathbb{E}_p\!\biggl[\sum_k d\alpha_k\,s_k
+\tfrac{1}{2}\sum_{kl}d\alpha_k d\alpha_l
\frac{\partial^2\log p}{\partial\alpha_k\partial\alpha_l}
\biggr]+\mathcal{O}(d\alpha^3).
\end{multline*}
The first-order term vanishes because \(\mathbb{E}_p[s_k]=0\).
The second-order term gives \(\tfrac{1}{2}d\bm{\alpha}^\top F\,d\bm{\alpha}\)
by the Hessian form. \(\square\)
Therefore \(F\) is the \textbf{Fisher--Rao metric} — the Riemannian
metric on the statistical manifold induced by the KL divergence.
\end{frame}

% ─────────────────────────────────────────────────────────────────────────────
\begin{frame}{Geometry of probability space}
Probability distributions form a \textbf{curved statistical manifold}.

\medskip

The score function \(\partial_\theta\log p(x|\theta)\) defines
the tangent vectors to this manifold: they measure how the
information content of an observation changes with model parameters.

\medskip

\begin{columns}[T]
\begin{column}{0.55\textwidth}
The \textbf{Fisher--Rao metric}
\[
ds^2=\sum_{ij}F_{ij}\,d\alpha_i\,d\alpha_j
\]
makes this manifold Riemannian.
Distances, curvatures, and information measures become
mathematically tractable in these coordinates.
\end{column}
\begin{column}{0.40\textwidth}
\textbf{Why logarithms?}

The logarithm transforms the positivity
constraint \(p\geq0\) and normalisation
\(\int p\,dx=1\) into an unconstrained,
geometrically flat coordinate system
where gradients measure relative changes.
\end{column}
\end{columns}
\end{frame}

% ── Quantum Fisher information ────────────────────────────────────────────────

\begin{frame}{Score for the VMC distribution}
In VMC the sampling distribution is
\(p(x;\bm{\alpha})=|\Psi(x;\bm{\alpha})|^2/Z(\bm{\alpha})\).

\medskip

\textbf{Lemma.}  Differentiating \(Z=\int|\Psi|^2\,dx\):
\[
\partial_{\alpha_k}\ln Z=2\langle O_k\rangle.
\]

Therefore the score for the VMC distribution is
\[
s_k(x;\bm{\alpha})
=\partial_{\alpha_k}\log p
=2\bigl(O_k(x)-\langle O_k\rangle\bigr)
\quad\text{(for real }\Psi\text{)}.
\]

The score is the \textbf{centred logarithmic derivative}
\(O_k=\partial_{\alpha_k}\ln\Psi\) — the same object that enters
the SR matrix through \(S_{ij}=\langle O_iO_j\rangle-\langle O_i\rangle\langle O_j\rangle\).
This is the key bridge between the SR and FIM formalisms.
\end{frame}

% ─────────────────────────────────────────────────────────────────────────────
\begin{frame}{Quantum Fisher information matrix equals \(4S\)}
Substituting the VMC score into the FIM definition:
\[
F_{ij}
=\mathbb{E}_p[s_i\,s_j]
=4\,\mathbb{E}_p\bigl[(O_i-\langle O_i\rangle)(O_j-\langle O_j\rangle)\bigr]
=4\bigl(\langle O_iO_j\rangle-\langle O_i\rangle\langle O_j\rangle\bigr).
\]

\[
\boxed{F_{ij}=4\,S_{ij}.}
\]

The factor of four arises from the squared coefficient in the score.
The SR matrix \(S\) is therefore \textbf{one quarter of the Fisher
information matrix} of the Born distribution \(p(x;\bm{\alpha})\).

\medskip

The Fubini--Study metric (natural distance on quantum state space)
gives the same result: for real \(\Psi\),
\(
ds^2_{\mathrm{FS}}=d\bm{\alpha}^\top S\,d\bm{\alpha},
\)
confirming that \(S\) governs both the quantum-state distance and
the KL distance between Born distributions.
\end{frame}

% ── Natural gradient ──────────────────────────────────────────────────────────

\begin{frame}{Natural gradient: derivation via constrained optimisation}
The \textbf{natural gradient} is the steepest-descent direction
in the Fisher--Rao metric: minimize \(\mathcal{L}\) subject to a
fixed KL step size,
\[
\min_{\delta\bm{\alpha}}\;\bm{g}^\top\delta\bm{\alpha}
\quad\text{s.t.}\quad
\delta\bm{\alpha}^\top F\,\delta\bm{\alpha}=\varepsilon^2.
\]

Lagrange multiplier \(\mu\) gives
\(\bm{g}-2\mu F\,\delta\bm{\alpha}=\bm{0}\), hence
\[
\boxed{\delta\bm{\alpha}_{\mathrm{NG}}=-\eta\,F^{-1}\bm{g}.}
\]

Plain gradient descent is the special case \(F=I\)
(Euclidean metric on parameter space).

\medskip

\textbf{Reparameterisation covariance.}
Under \(\bm{\beta}=f(\bm{\alpha})\) with Jacobian \(J\),
the FIM transforms as \(F^{(\beta)}=J^\top F^{(\alpha)}J\),
so \(\delta\bm{\beta}_{\mathrm{NG}}=J^{-1}\delta\bm{\alpha}_{\mathrm{NG}}\):
the natural-gradient update is \emph{covariant}.
Plain gradient descent is not.
\end{frame}

% ─────────────────────────────────────────────────────────────────────────────
\begin{frame}{Three-way equivalence}
Since \(F_{ij}=4S_{ij}\), the natural gradient update
\(\delta\bm{\alpha}_{\mathrm{NG}}=-\eta F^{-1}\bm{g}=-\frac{\eta}{4}S^{-1}\bm{g}\)
recovers the SR update when \(\eta/4=\delta\tau\).
All three descriptions are the \emph{same} update:

\medskip

\begin{center}
\renewcommand{\arraystretch}{1.8}
\begin{tabular}{c@{\quad}c@{\quad}c}
\textbf{Projected imaginary time} & & \textbf{Natural gradient (statistics)} \\
\(S\,\delta\bm{\alpha}=-\delta\tau\,\bm{g}\)
& \(\boldsymbol{\Longleftrightarrow}\) &
\(\delta\bm{\alpha}=-\dfrac{\delta\tau}{4}\,F^{-1}\bm{g}\) \\[6pt]
\multicolumn{3}{c}{\(\boldsymbol{\Updownarrow}\)} \\[2pt]
\multicolumn{3}{c}{\textbf{Natural gradient (quantum geometry)}} \\
\multicolumn{3}{c}{\(\delta\bm{\alpha}=-\delta\tau\,S^{-1}\bm{g}\)}
\end{tabular}
\end{center}

\smallskip

Quantum mechanics (ITE projection), information geometry (steepest
descent on the statistical manifold), and optimisation theory
(natural gradient) all say the same thing.
\end{frame}

% ─────────────────────────────────────────────────────────────────────────────
\begin{frame}{Cram\'er--Rao bound and principled regularization}
The FIM controls \textbf{how precisely parameters can be estimated}.

\medskip

\textbf{Cram\'er--Rao bound:} for any unbiased estimator
\(\hat{\bm{\alpha}}\) from \(N_s\) i.i.d.\ samples,
\[
\mathrm{Cov}(\hat{\bm{\alpha}})\succeq\frac{1}{N_s}F(\bm{\alpha})^{-1}.
\]

In VMC: directions with large \(F_{kk}\) can be estimated precisely
and updated aggressively; directions with small \(F_{kk}\) cannot be
estimated beyond the CR limit and should be damped.

\medskip

This gives a \textbf{principled interpretation of regularization}:
\[
(S+\lambda I)^{-1}=\bigl(\tfrac{1}{4}F+\lambda I\bigr)^{-1}
\]
damps parameter directions below the Cram\'er--Rao estimation floor,
interpolating between SR (\(\lambda\to 0\)) and plain gradient
descent (\(\lambda\to\infty\)).
\end{frame}

% ─────────────────────────────────────────────────────────────────────────────
\begin{frame}{Geometry and connections summary}
\begin{center}
\small
\setlength{\tabcolsep}{5pt}
\begin{tabular}{@{}lll@{}}
\toprule
\textbf{Object} & \textbf{Expression} & \textbf{Geometric / statistical role} \\
\midrule
Score function
  & \(s_k=\partial_{\alpha_k}\log p\)
  & relative sensitivity; tangent to stat.\ manifold \\
Information content
  & \(I(x)=-\log p(x)\)
  & additive; large for rare events \\
Shannon entropy
  & \(S=-\mathbb{E}[\log p]\)
  & expected information content \\
Fisher information matrix
  & \(F_{ij}=\mathbb{E}[s_is_j]\)
  & curvature of log-likelihood \\
Fisher--Rao metric
  & \(ds^2=d\bm{\alpha}^\top F\,d\bm{\alpha}\)
  & distance on statistical manifold \\
KL divergence (2nd order)
  & \(\dkl\approx\tfrac{1}{2}d\bm{\alpha}^\top F\,d\bm{\alpha}\)
  & KL curvature \(\equiv\) Fisher metric \\
VMC score
  & \(s_k=2(O_k-\langle O_k\rangle)\)
  & centred log-derivative \\
SR matrix vs FIM
  & \(F_{ij}=4S_{ij}\)
  & quantum Fisher = 4\(\times\) SR metric \\
Fubini--Study metric
  & \(ds^2_{\mathrm{FS}}=d\bm{\alpha}^\top S\,d\bm{\alpha}\)
  & quantum state distance \\
Natural gradient
  & \(\delta\bm{\alpha}=-\eta F^{-1}\bm{g}\)
  & steepest descent on stat.\ manifold \\
SR update
  & \(S\,\delta\bm{\alpha}=-\delta\tau\,\bm{g}\)
  & projected ITE \(\equiv\) natural gradient \\
Cram\'er--Rao bound
  & \(\mathrm{Cov}(\hat{\bm{\alpha}})\succeq(N_sF)^{-1}\)
  & estimation precision limit \\
\bottomrule
\end{tabular}
\end{center}
\end{frame}

% ══════════════════════════════════════════════════════════════════════════════
%  PART III  –  REGULARIZATION AND ALGORITHM
% ══════════════════════════════════════════════════════════════════════════════

\begin{frame}{Regularization and practical solution}
In practice \(S\) may be singular or ill-conditioned because it is
estimated from a finite Monte Carlo sample.

A standard fix is diagonal regularization:
\[
S\;\to\;S+\lambda I,\qquad\lambda>0.
\]

Then solve
\[
(S+\lambda I)\,\delta\bm{\alpha}=-\delta\tau\,\bm{g}.
\]

Interpretation:
\begin{itemize}
  \item stabilizes inversion,
  \item damps directions with negligible Fisher information
        (flat directions of the statistical manifold),
  \item interpolates between SR (\(\lambda\to 0\)) and ordinary
        gradient descent (\(\lambda\to\infty\)).
\end{itemize}
\end{frame}

% ─────────────────────────────────────────────────────────────────────────────
\begin{frame}{Algorithmic summary}
At each iteration:
\begin{enumerate}
  \item Sample configurations \(x^{(\ell)}\) from
        \(P(x)=|\Psi(x)|^2/\braket{\Psi}{\Psi}\).
  \item Compute logarithmic derivatives
        \(O_k(x^{(\ell)})=\partial_{\alpha_k}\ln\Psi(x^{(\ell)})\).
  \item Compute local energies
        \(\Eloc(x^{(\ell)})=(H\Psi)(x^{(\ell)})/\Psi(x^{(\ell)})\).
  \item Estimate
  \[
  S_{ij}=\langle O_iO_j\rangle-\langle O_i\rangle\langle O_j\rangle,
  \qquad
  g_i=\langle O_i\Eloc\rangle-\langle O_i\rangle\langle\Eloc\rangle.
  \]
  \item Solve \((S+\lambda I)\,\delta\bm{\alpha}=-\delta\tau\,\bm{g}\).
  \item Update \(\bm{\alpha}\leftarrow\bm{\alpha}+\delta\bm{\alpha}\).
\end{enumerate}
\end{frame}

% ─────────────────────────────────────────────────────────────────────────────
\begin{frame}{Summary}
\begin{columns}[T]
\begin{column}{0.48\textwidth}
\textbf{Stochastic reconfiguration}
\begin{itemize}\setlength\itemsep{3pt}
  \item ITE selects the ground state;
        projecting onto \(\mathcal{M}\) gives
        \(S\,\delta\bm{\alpha}=-\delta\tau\,\bm{g}\).
  \item \(S_{ij}\): covariance of log-derivatives.
  \item \(g_i\): covariance with the local energy.
\end{itemize}

\textbf{Why \(\log p(x)\) is fundamental}
\begin{itemize}\setlength\itemsep{3pt}
  \item Converts products into sums.
  \item Measures relative sensitivity \(\delta p/p\).
  \item Makes information additive.
  \item Connects energy to probability (stat.\ mech.).
  \item Defines the natural geometry on
        probability space.
\end{itemize}
\end{column}
\begin{column}{0.48\textwidth}
\textbf{Fisher information and natural gradient}
\begin{itemize}\setlength\itemsep{3pt}
  \item Score \(s_k=\partial_{\alpha_k}\log p\)
        has zero mean; its covariance is the FIM.
  \item FIM = curvature of KL divergence
        = Fisher--Rao metric.
  \item For VMC: \(F_{ij}=4S_{ij}\);
        SR matrix \(S\) is \(\tfrac{1}{4}\) of the FIM.
  \item SR update = natural gradient
        = projected imaginary-time evolution.
  \item Natural gradient is reparameterisation-covariant;
        GD is not.
  \item Regularization \((S{+}\lambda I)^{-1}\) damps
        directions below the Cram\'er--Rao floor.
\end{itemize}
\end{column}
\end{columns}
\end{frame}

\end{document}
